% Full English translation of the consolidated continuation proofs.
\section{The minimum slice degree theorem}\label{sec:jets}
Set $V=\gamma(u)-t^\mu X$. For every $P\circ\gamma=u$, there is a polynomial slice
\[
S_{P,\mu}=\frac{u-P(V)}{t^\mu},\qquad D_\mu(S_{P,\mu})=1.
\]
Polynomiality follows from exact divisibility of the numerator; the displayed denominator does not make this a nonpolynomial rational function. This classical construction goes back to van den Essen--van Rossum \cite{vdER}. The calculations below refine the degree-minimization problem.

\subsection*{Theorem 3.1. Completeness of the reduction}
The minimum degree among all polynomial slices equals the minimum of $\deg S_{P,\mu}$ over all parameter polynomials $P$.

\textit{Proof.} Introduce the weights
$\wt(t)=1$, $\wt(u)=0$, $\wt(X_i)=-\mu$. The derivation raises weight by $\mu$. If $D_\mu(s)=1$, the component $s_{-\mu}$ of weight $-\mu$ is also a slice; its ordinary degree is no greater than that of $s$.
After inverting $t$,
\[
A_t=k[t,t^{-1},u,V],\qquad D_\mu=t^\mu\partial_u.
\]
The element $t^\mu s_{-\mu}-u$ is invariant and has weight zero. Since all $V_i$ have weight zero, it equals $-P(V)$ for a unique $P\in k[V]$. Reducing the exact polynomial identity modulo $t$ gives $P\circ\gamma=u$.\hfill$\square$

\subsection*{Theorem 3.2. Exact formula and finite criterion}
With the convention $\deg0=-\infty$, we have
\[
\boxed{\deg S_{P,\mu}=\max_{\alpha\ne0}
\bigl\{\deg_u((\partial^\alpha P)\circ\gamma)
+(\mu+1)|\alpha|-\mu\bigr\}.}
\]
Indeed, Taylor's formula gives
\[
S_{P,\mu}=\sum_{\alpha\ne0}\frac{(-1)^{|\alpha|+1}}{\alpha!}
t^{\mu(|\alpha|-1)}X^\alpha(\partial^\alpha P)(\gamma(u)).
\]
Distinct $\alpha$ correspond to distinct monomials in $X$, so they cannot cancel one another. A nonzero derivative of order $\deg P=r$ gives the lower bound $(\mu+1)r-\mu$.

Consequently, a slice of degree $\le d$ exists if and only if the following finite linear system is solvable:
\[
\begin{gathered}
\deg P\le\floor{\frac{d+\mu}{\mu+1}},\qquad P\circ\gamma=u,\\
[u^j](\partial^\alpha P)(\gamma(u))=0\quad
\text{for }j>d-(\mu+1)|\alpha|+\mu,\quad\alpha\ne0.
\end{gathered}
\]
It is linear in the coefficients of $P$. For $n=3$, the number of unknowns is at most
$\binom{\lfloor(d+\mu)/(\mu+1)\rfloor+3}{3}$.
This is a complete criterion, rather than a search restricted to a chosen form of slice.

\subsection*{Theorem 3.3. All sufficiently large multiplicities}
Let
\[
M=\max_i\deg f_i,\qquad
r_\gamma=\min\{\deg P:P\circ\gamma=u\},\qquad
m_\gamma(\mu)=\min_{D_\mu(s)=1}\deg s.
\]
Then
\[
\boxed{m_\gamma(\mu)=(\mu+1)r_\gamma-\mu\quad\text{for }\mu\ge M-1.}
\]
For a nonconstant $P$ of degree $r$, a derivative of order $q$ has degree at most $M(r-q)$ after substitution. When $\mu+1\ge M$,
\[
M(r-q)+(\mu+1)q-\mu\le(\mu+1)r-\mu.
\]
A nonzero derivative of order $r$ attains this bound. Minimizing over $P$ completes the proof. Moreover, $m_\gamma(\mu)$ is nondecreasing: every expression in the maximum has nonnegative slope $|\alpha|-1$.

This degree measure depends on the original coordinates. It is not an invariant of conjugacy of LNDs and, without a further argument, cannot serve as an obstruction to polynomiality of the kernel.

\section{Exact calculations for harmonic models}\label{sec:computed}
For $\mu=1$, the following results were obtained. The lower bounds are certified by rational dual certificates, and the upper bounds by explicit slices.

\begin{center}\small\renewcommand{\arraystretch}{1.22}
\begin{tabular}{@{}lrrr@{}}\toprule
Model & $m_\gamma(1)$ & Unknowns & Nonzero weights\\\midrule
$H(3,4,5)$ &7&20&11\\
$H(3,5,7)$ &11&56&35\\
$H(4,5,7)$ &12&84&41\\
$H(3,7,8)$ &13&84&41\\
$H(3,10,11)$ &19&220&111\\
$H(3,13,14)$ &25&455&225\\
$H(5,7,8)$ &11&56&29\\
$H(7,10,11)$ &15&120&61\\\bottomrule
\end{tabular}\end{center}
The last two columns refer to the system excluding degree one below the minimum. The first two rows reproduce earlier results with more compact certificates.

If the system is written as $A c=b$, a lower-bound certificate is a rational vector $\lambda$ satisfying $\lambda^{\mathsf T}A=0$ and $\lambda^{\mathsf T}b=1$. This rules out a solution over every field of characteristic zero, not only over $\Q$. A separate verification script rebuilds the columns and performs the calculations using \file{Fraction}; no symbolic solver is needed for verification.

\subsection*{$H(3,7,8)$: a lower-degree retraction is not always better}
The Chebyshev identities give
\[
C_7C_8=C_{15}+C_1,\qquad C_5(C_3)=C_{15},\qquad
P_5=yz-C_5(x)=yz-x^5+5x^3-5x.
\]
The resulting slice of degree $13$ is optimal for $\mu=1$. However, the minimum degree of the parameter polynomial itself is $4$; one may take
\[
\begin{split}
P_4={}&x^3z+x^3-x^2yz-xy^2-2xz-x\\
&+y^3+2yz-2y.
\end{split}
\]
For $\mu=1$, this $P_4$ yields a slice of degree $19$. Thus minimizing $\deg P$ and minimizing $\deg S_P$ are different problems.

The full dependence on multiplicity is
\[
\boxed{\begin{aligned}
m_{378}(1)&=m_{378}(2)=13,\quad m_{378}(3)=17,\quad m_{378}(4)=21,\\
m_{378}(5)&=22,\quad m_{378}(6)=23,\quad
m_{378}(\mu)=3\mu+4\quad(\mu\ge7).
\end{aligned}}
\]
For $\mu\le4$, $P_5$ suffices; from $\mu=5$ onward, $P_4$ suffices. The lower bounds for $1\le\mu\le7$ have been checked exactly; the last test also excludes $\deg P\le3$. Theorem 3.3 applies thereafter.

\subsection*{$H(4,5,7)$: a formula for all multiplicities}
Here
\[
P_3=(x^2-2)z-y^3+3y,\qquad
\boxed{m_{457}(\mu)=\max\{12,\mu+9,2\mu+3\}.}
\]
The first, second, and third derivatives of $P_3$ along the curve have maximum degrees $11$, $7$, and $0$. This gives the upper bound. For $\mu=1,2,3$, the lower bound $12$ follows from the certificate at $\mu=1$ and monotonicity.

Restriction to the curve is injective on polynomials of degree $\le3$: the corresponding $20\times20$ minor has determinant $267386$. Therefore $P_3$ is unique in this degree range. For $\mu=4,5$, every $P$ of degree $\ge4$ gives a slice of degree at least $3\mu+4>\mu+9$. For $\mu\ge6$, the general theorem applies with $M=7$, $r_\gamma=3$.

\subsection*{An infinite family of upper bounds}
For odd $a,m\ge3$, set $b=(ma-1)/2$, $c=(ma+1)/2$. The indices are pairwise coprime, and $P=yz-C_m(x)$ recovers the parameter. The exact degree of its slice is
\[
\deg S_{P,\mu}=\max\{a(m-1)+1,(\mu+1)m-\mu\}.
\]
For $(a,b,c)=(3,3n+1,3n+2)$, this gives $m_\gamma(1)\le6n+1$. Equality is certified for $n=1,2,3,4$; the general lower bound has not yet been proved.

\section{The kernel and formal triviality}\label{sec:formal}
Let $I=\ker(k[Y_1,\ldots,Y_n]\to k[u])$ be the ideal of the embedding, and let $B_\mu=\ker D_\mu$. Taylor expansion yields the following description, classical in nature:
\[
\boxed{B_\mu=k[t,V,I(V)/t^\mu]\simeq k[Y][t,I/t^\mu].}
\]
For $\mu=1$, the corresponding construction is described in \cite[Proposition 5.4, Corollary 5.5]{vdER}. It does not assert polynomiality of $B_\mu$.

The elementary key is this: if all derivatives of $F$ of orders $<q$ vanish along $\gamma$, then $F\in I^q$. To prove this, choose $P\circ\gamma=u$ and put $g_i=Y_i-f_i(P(Y))$. Taylor expansion about $\gamma(P(Y))$ first shows that $I=(g_1,\ldots,g_n)$, and then gives the required membership. The converse follows from the Leibniz rule. For $j<0$, a weight component $t^jF(V)$ is polynomial exactly when $F\in I^{\lceil-j/\mu\rceil}$; this gives the kernel formula.

\subsection*{Theorem 5.1. All finite thickenings are trivial}
For every embedding, every $\mu\ge1$, and every $N\ge1$,
\[
\boxed{B_\mu/(t^N)\simeq(k[t]/(t^N))^{[n]}.}
\]
The isomorphisms can be chosen compatibly in $N$, so
\[
\boxed{\widehat B_\mu^{\,(t)}\simeq k[p,q_1,\ldots,q_{n-1}][[t]].}
\]
The right-hand side consists of power series in $t$ with polynomial coefficients. The degrees of these coefficients in $p,q$ may grow. This is not a claim about the ordinary polynomial ring $k[[t]][p,q]$.

\textit{Proof.} Set $a(u)=\nabla P(\gamma(u))$, $v(u)=\gamma'(u)$. The chain rule gives $a v=1$. Over the PID $k[u]$, the module of rows annihilating $v$ is free of rank $n-1$; choose a basis $\ell_1,\ldots,\ell_{n-1}$. Together with $a$, these rows form an invertible matrix over $k[u]$.

A row $\ell$ can be lifted explicitly to the ideal of the curve:
\[
F_\ell(Y)=\sum_i\ell_i(P(Y))\bigl(Y_i-f_i(P(Y))\bigr).
\]
We have $F_\ell\in I$ and $\nabla F_\ell(\gamma)=\ell-(\ell v)a=\ell$. Define
\[
p=P(V),\qquad q_j=-F_{\ell_j}(V)/t^\mu,\qquad s=(u-P(V))/t^\mu.
\]
These are polynomials; $p,q_j$ are invariant, and $D_\mu(s)=1$. At $t=0$, the change $(u,X)\mapsto(p,q,s)$ is
\[
(u,X)\longmapsto(u,\ell_1(u)X,\ldots,\ell_{n-1}(u)X,a(u)X)
\]
and is invertible.

After composing with its polynomial inverse at $t=0$, we obtain an endomorphism $\Phi\equiv\mathrm{id}\pmod{t^\mu}$. On $A/(t^N)$, the operator $E=\mathrm{id}-\Phi$ raises $t$-adic order by $\mu$, so $E^L=0$ for $L=\lceil N/\mu\rceil$. The inverse map is the finite sum $\sum_{j=0}^{L-1}E^j$. It is the inverse of a ring homomorphism as an additive map and is therefore automatically multiplicative.

Thus
\[
A/(t^N)=(k[t]/(t^N))[p,q_1,\ldots,q_{n-1},s],\qquad\overline D_\mu=\partial_s.
\]
The global slice theorem gives $A=B_\mu[s]$. Consequently, the kernel of the reduced derivation is $B_\mu/(t^N)$. The same polynomials $p,q$ work for every $N$, ensuring compatibility.\hfill$\square$

\textbf{Consequence for the strategy.} The abstract type of any one finite thickening of the fiber $t=0$ cannot distinguish these kernels from a polynomial kernel. One must control global algebraization, degrees of the inverse map, or additional structures. Finite certificates for degree restrictions remain useful: they retain the original coordinates, which are absent from the abstract thickening ring.

\section{A restriction on coordinates for the cubic family}\label{sec:uniform}
Let $\gamma_n=(C_3,C_{3n+1},C_{3n+2})$, $n\ge1$. Over $\C$, there is no polynomial coordinate $H$ satisfying $H\circ\gamma_n=u$ of either form
\[
\boxed{H=Q(x,y)+M(x)z\quad\text{or}\quad H=Q(x,z)+M(x)y.}
\]
The degrees of $Q$ and $M$ are unrestricted. The exclusion also applies to an absolute coordinate: the full automorphism is not required to preserve $x$.

\subsection*{Proof of the relative statement}
Write $U_j=U_j(x/2)$ for the Chebyshev polynomials of the second kind. In the free module with basis $1,u,u^2-2$ over $\C[x]$, where $u^3-3u=x$, we have
\[
y=U_nu-U_{n-1}(u^2-2),\qquad
z=-U_{n-1}u+U_n(u^2-2).
\]
For the first form, reduce $Q$ modulo the monic relation
$y^3-3y-C_{3n+1}(x)=0$ to $Ay^2+By+C$. Comparing coefficients gives
\[
U_{3n}A+U_{2n}M=U_{n-1}.\tag{6.1}
\]
When $M=0$, this equality is impossible by degree. If $(x,H,W)$ is an automorphism and $M(\alpha)=0$, the polynomial $Q(\alpha,y)$ must be affine-linear in $y$. Hence $A(\alpha)=0$ and $U_{n-1}(\alpha)=0$.

All roots of $M$ belong to the real roots of $U_{n-1}$. The two largest roots of $U_{3n}$,
$x_j=2\cos(j\pi/(3n+1))$, $j=1,2$, lie to the right of all these roots. For $n=1$, $U_0$ has no roots at all. Substitution in (6.1), together with the sine formula, gives
$M(x_1)=1$, $M(x_2)=-1$.
But for $M=c\prod(x-\alpha)^{e_\alpha}$, the ratio $M(x_1)/M(x_2)$ is positive, even when $c$ is complex. This is a contradiction.

For the second form, one similarly obtains
\[
U_{3n+1}A+U_{2n}M=U_n.
\]
The two largest roots of $U_{3n+1}$ lie to the right of the roots of $U_n$ and again force opposite values of $M$. This completes the relative proof.

\subsection*{Passage to an absolute coordinate}
For $M\ne0$, the polynomial $H$ is a coordinate over $\C(x)$. If it were an absolute coordinate, Proposition 3.1 of Edo--van den Essen--Maubach \cite{EVM} would make it a coordinate over $\C[x]$. This contradicts what has just been proved.

For $n=1$, this exclusion is compatible with the project's complex constructions for $H(3,4,5)$: the parameter component must depend nonlinearly on both transverse variables. This restricts the form of a rectifier; it does not prove that no rectifier exists.

\subsection*{The natural modification also fails}
Consider
\[
\beta_n=(x,C_n(x)z-y^2+2,y),\qquad
\beta_n\circ\gamma_n=(C_3,C_2,C_{3n+1}).
\]
At every root $\alpha$ of $C_n$, the center $H(\alpha,2-y^2,y)$ cannot be affine-linear in $y$ if $H$ recovers the parameter of the target triple. Otherwise, in the quotient by $u^3-3u-\alpha$, we obtain
\[
u=\lambda\bigl(U_n(\alpha/2)u-U_{n-1}(\alpha/2)(u^2-2)\bigr)+\nu.
\]
Here $U_{n-1}(\alpha/2)\ne0$, because
$C_n(x)^2+(4-x^2)U_{n-1}(x/2)^2=4$. Comparing coefficients first gives $\lambda=0$ and then a contradiction. This proof imposes no bound on the degree of $H$.

For odd $n$, the target triple is not an embedding: $u=\pm\sqrt3$ have the same image. For even $n=2m$, the target has an explicit relative coordinate. If
$C_{6m+1}(u)=uP_m(u^2-2)$, $\varepsilon=P_m(1)=(-1)^m$, then
\[
A_m(Y)=\frac{1-\varepsilon P_m(Y)}{Y-1},\qquad
H_0=xA_m(Y)+\varepsilon Z
\]
recovers the parameter. Thus even rectifiability of the target does not repair the bad center of this particular transfer.

\subsection*{An additional finite result for $H(3,7,8)$}
Coordinates $Q(x,L)+M(x)N$ are excluded when $(L,N)$ is any invertible constant linear change of $(y,z)$, $\deg M\le10$, and the degree of $Q$ is arbitrary. After normalizing $L=y+az$, the problem reduces to 56 exact conditions on the multiplicities of the roots of $M$. In every case, the gcd of the necessary polynomial conditions is $1$; repeated roots at $a=\pm i/2$ are checked separately. This is a \textbf{bounded result}: the case $\deg M>10$ is not excluded here.

\section{H(4,5,7): the degree-four barrier in the relative class}\label{sec:h457}
Suppose that $F\circ(C_4,C_5,C_7)=u$ and $\deg F\le4$. Then:
\begin{itemize}
\item the unique such $F$ that is a coordinate of $k(x)[y,z]$ is $F_0=(x^2-2)z-y^3+3y$;
\item no such $F$ is a coordinate of $k(y)[x,z]$ or $k(z)[x,y]$;
\item $F_0$ is not a coordinate over $k[x]$: its derivatives with respect to $y,z$ vanish simultaneously at $x^2=2$, $y^2=1$.
\end{itemize}
Consequently, any method that completes one of the pairs $(x,F)$, $(y,F)$, or $(z,F)$ to an automorphism necessarily requires $\deg F\ge5$.

The proof exhausts the entire space of polynomials of degree $\le4$. The restriction matrix for the 35 monomials has rank 28. All parameter polynomials form the affine space $F_0+\langle K_0,\ldots,K_6\rangle$. The archive lists all seven $K_i$ explicitly; completeness is verified by nonzero minors.

The Jung--van der Kulk theorem is then applied: a plane coordinate of degree 3 has the form $P_3(L)+cM$, and one of degree 4 has either the form $P_4(L)+cM$ or $A(M+Q_2(L))^2+bL+c$, with the appropriate nonzero coefficients. Successive coefficient comparisons in all three projections exclude the possibilities. In the final branch, the nonzero resultant $6561301239$ is computed; an explicit B\'ezout identity with constant $729033471$ is also supplied.

The verifier checks 35 exact conditions, and a separate audit repeats the ranks, boundary cases, and parameter eliminations. The structural source of the normal forms is the classical theorem and its degree-divisibility lemma; a convenient primary account is given in \cite{Chau}.

\textbf{Scope of the assertion.} It does not exclude an absolute coordinate of degree 4 whose completion changes all three original coordinates, nor does it exclude relative coordinates of degree $\ge5$.

\section{The complete filtration of a one-step transfer}\label{sec:transfer}
Let $R=k[x]$, $A=R[y,z]$, $(H,W)\in\Aut_R R[Y,Z]$, and
\[
d=H_YW_Z-H_ZW_Y\in k^*,\qquad
\beta=(p(y)+mz,y),\quad p\in R[y],\ 0\ne m\in R.
\]
Set $h=H\circ\beta$, $w=W\circ\beta$, and $D=h_z\partial_y-h_y\partial_z$.
Extract the content of the nonconstant monomials:
\[
c=h(x,0,0),\qquad h=c+g h_0,
\]
where $g\in R$ is the monic gcd of the coefficients of all monomials of $h$ that are nonconstant in $y,z$.

\subsection*{Theorem 8.1}
We have $g\mid m$ and $B=\ker D=k[x,h_0]$. For every irreducible $\pi\mid m_1=m/g$, define
\[
E_\pi=v_\pi(m_1),\qquad e_\pi=\deg_y(h_0\bmod\pi)\ge1.
\]
The reduction of $h_0$ here is independent of $z$. For $\mathcal F_n=\ker D^{n+1}$,
\[
\boxed{D^n(\mathcal F_n)=
\left(g^n\prod_{\pi\mid m_1}\pi^{E_\pi(n-\lfloor n/e_\pi\rfloor)}\right)B.}
\]
In particular,
\[
\boxed{\pl(D)=D(A)\cap B=
\left(g\prod_{e_\pi\ge2}\pi^{E_\pi}\right)B.}
\]
All multiplicities of the exceptional divisor are retained. Good components with $e_\pi=1$ disappear from the plinth ideal of the primitive derivation; the factor $g$ is accounted for separately.

\subsection*{Main steps of the proof}
\textit{1. Kernel and normalization.} Over $k(x)$, the pair $(h,w)$ is a coordinate system, and $D(w)=dm$. It follows that $g\mid m$, and the kernel over this field is $k(x)[h_0]$. The reduction of $h_0$ modulo every prime $\pi\in R$ is nonconstant; hence $R[h_0]$ is $\pi$-saturated in $A$. Successive cancellation of denominators gives $A\cap k(x)[h_0]=R[h_0]$.

\textit{2. An exact equation for the modification.} Let $(\mathcal Y,\mathcal Z)$ be the inverse of $(H,W)$. The polynomial
\[
N(T,V)=\frac{\mathcal Y(c+gT,V)-p(\mathcal Z(c+gT,V))}{g}
\]
does belong to $R[T,V]$, and $\gcd(N,m_1)=1$; this is proved by reduction modulo $g$ and invertibility of the center over the finite quotient ring. Then
\[
A\simeq R[T,V,z]/(m_1z-N(T,V)),\qquad T=h_0,\quad V=w.
\]
Polynomiality of $N$ is essential: formal division without checking it would leave a gap.

\textit{3. A monic model at each prime.} Localize only $R$: let $O=R_{(\pi)}$, with coefficient ring $O[T]$. Over $O/(\pi^{E_\pi})$, the center has the form $(G(y),J(y))$ and generates the entire polynomial ring in $y$. This ring is free of rank $e_\pi$ over the ring in $G$; the proof repeats ordinary division and uses a finite lifting process through powers of $\pi$.
The characteristic polynomial of multiplication by $J$ gives a monic $P_\pi(T,V)$ of degree $e_\pi$ in $V$. The equations of the center and $P_\pi$ generate the same ideal modulo $\pi^{E_\pi}$, leading to the exact model
\[
P_\pi(h_0,w)=\pi^{E_\pi}v.
\]
Conditions depending on $T$ are retained here: replacing $O[T]$ by its rational function field in $T$ would not suffice.

\textit{4. Basis and leading coefficient.} The monomials $w^iv^j$, $0\le i<e_\pi$, form a basis over $O[h_0]$. Their $D$-degrees are $i+e_\pi j$ and are pairwise distinct. Thus leading terms cannot cancel. When $n=i+e_\pi j$, we obtain exponent $E_\pi(n-j)$. Intersecting the local conditions over all primes of $R$, and using the Chinese remainder theorem to construct preimages, gives the global formula.

The full algebraic proof, including divisibility, monic lifting, and gluing, is provided in \file{relative_lnd/PROOF.md}; a separate audit is in \file{literature/independent_image_ideal_review.md}. In addition, 16 examples and 28 separate exact conditions were checked.

\subsection*{An example with a repeated divisor and a nonlinear lift}
Let
\[
h=y^2+xy^3+x^2z,\qquad D=x^2\partial_y-(2y+3xy^2)\partial_z.
\]
Here $E=2$ and $e=2$. For $v=y^4+xyz-z$,
\[
Dv=2y+xh,\qquad D^2v=2x^2,\qquad
z=-(1+xy)v+hy^2.
\]
The sequence of image ideals has exponents $0,2,2,4,4,6,6,\ldots$. Finding a single image element yields only an inclusion; equality of the entire ideal follows from the theorem.

\subsection*{Consequences for the theory of LNDs}
For $D_{M,e}=x^M\partial_y-e y^{e-1}\partial_z$, $M\ge1$, $e\ge2$,
\[
I_n(D_{M,e})=(x^{M(n-\lfloor n/e\rfloor)}).
\]
The exponent $M$ is recovered from the multiplicity of the prime divisor of the plinth ideal, and $e$ is recovered as the first $n$ for which $I_n\supsetneq I_1^n$. Thus distinct pairs $(M,e)$ give nonconjugate LNDs. The family itself is classical; no priority is claimed for this formulation of the consequence.

For $h=(x^2-2)z-y^3+3y$, which arises for $H(4,5,7)$,
\[
\ker D=k[x,h],\qquad I_n(D)=(x^2-2)^{n-\lfloor n/3\rfloor}k[x,h].
\]
This is exact information about this particular LND, not an obstruction to every possible rectifier of the curve.

\section{Iterated modifications and Cremona groups}\label{sec:iterated}
Define
\[
r_{-1}=z,\quad r_0=y,\quad r_j=xr_{j-2}-r_{j-1}^2,\qquad j\ge1.
\]
For a fixed $\ell\ge1$, set $p=r_\ell$, $w=r_{\ell-1}$, and
\[
D_\ell=(-1)^\ell\bigl(p_y\partial_z-p_z\partial_y\bigr).
\]
Then $D_\ell(w)=x^\ell$ and $\ker D_\ell=B=k[x,p]$. The derivation is irreducible: a common divisor of its components must divide $x^\ell$, while $p\bmod x=-y^{2^\ell}$ has nonzero derivative with respect to $y$.

On the open set $x\ne0$, every operation is invertible; denominators appear when extending across $x=0$. This gives a concrete connection with Cremona transformations: birational invertibility is straightforward, whereas polynomiality and factorization length are controlled by the exceptional divisor.

\subsection*{Theorem 9.1. The entire filtration}
For $I_n(D)=D^n(\ker D^{n+1})$,
\[
\boxed{I_n(D_\ell)=x^{a_\ell(n)}B,\qquad
a_\ell(n)=\ell n-\sum_{j=1}^{\ell}\floor{\frac n{2^j}}.}
\]
The proof gives an explicit basis. Set $v_j=r_{\ell-1-j}$ for $0\le j\le\ell$, and $v_{-1}=p$. Then
\[
v_j^2=xv_{j+1}-v_{j-1}\quad(0\le j<\ell),\qquad \deg_{D_\ell}v_j=2^j.
\]
By the monic relations, the normal monomials have the form
\[
s_n=\prod_{j=0}^{\ell-1}v_j^{\varepsilon_j}\,v_\ell^q,
\qquad n=\sum_{j=0}^{\ell-1}\varepsilon_j2^j+q2^\ell,
\quad\varepsilon_j\in\{0,1\}.
\]
They form a $B$-basis and have pairwise distinct $D$-degrees. When expressed in terms of $w$, the leading coefficient of $v_j$ is $x^{-(2^j-1)}$. Applying $D_\ell^n$ and using the binary expansion of $n$ gives the stated sum of floor functions; the factor $n!$ is a unit in $k$.

For $\ell=3$, the exponents for $0\le n\le8$ are
\[
0,\ 3,\ 5,\ 8,\ 9,\ 12,\ 14,\ 17,\ 17.
\]
For $\ell\ge2$, the derivation $D_\ell$ is not conjugate to any irreducible one-step transfer of Theorem 8.1, even after changing the base coordinate. Indeed, the unique prime divisor of the plinth ideal and its multiplicity $\ell$ are preserved by every conjugacy. For the one-step class, the exponent of $I_2$ at this prime would be $\ell$ when $e=2$, or $2\ell$ when $e>2$, whereas here it equals $2\ell-1$.

In particular, $D_\ell$ is \textbf{not triangularizable} for $\ell\ge2$. If, in some coordinates, it had the form $a(X)\partial_Y+b(X,Y)\partial_Z$ with $a\ne0$, it would belong to the one-step class with $h=a(X)Z-\int b(X,Y)\,dY$. If the derivative of the first coordinate is a nonzero constant, the plinth ideal equals the entire kernel; when $a=0$, irreducibility likewise gives a constant translation. Both cases are incompatible with our plinth ideal. The existence of such nontriangularizable examples is not claimed as new; the point here is an explicit way to detect them using two ideals.

\subsection*{The associated graded ring}
The same monic relations give a fuller description:
\[
\boxed{\gr_{D_\ell}A\simeq
B[\xi_0,\ldots,\xi_\ell]/(\xi_j^2-x\xi_{j+1}:j<\ell),\quad
\deg\xi_j=2^j.}
\]
After reduction modulo $x$, the nilradical $N=(\xi_0,\ldots,\xi_{\ell-1})$ satisfies $N^\ell\ne0$, $N^{\ell+1}=0$, and $N/N^2$ is free of rank $\ell$ over $k[p,\xi_\ell]$.

An illustrative example is given by $D_3$ and $xD_2$: they have the same plinth ideal $(x^3)$, and the exponents of $I_n$ agree for $n\le7$. At $n=8$, they are $17$ and $18$, respectively. The full filtration detects a difference that the first seven ideals do not.

This illustrates the additional information in the filtration, rather than a newly solved difficult conjugacy problem: the fields are already distinguished by irreducibility.
For $\ell=2$, nontriangularizability is also known from
\cite[Example 4.19]{DG}; the novelty of the general family and formulas is not established here.

\subsection*{Connection with the current question on complexity}
In the journal article by Dasgupta--Gaifullin \cite[Theorem 4.12]{DG}, universal level is related to the complexity of a birational representation by $\operatorname{ulev}(D)=c(\alpha)+1$. Their Question 4.21 concerns a candidate constructed by three alternating modifications and permutations.

Our recurrence allows this type of construction to be studied through its exact filtration. \textbf{However, the equality $c(\varphi)=3$ is not proved here}: it requires excluding every shorter admissible factorization, not merely the canonical one-step class. Extracting such a lower bound from $\gr_DA$ is a substantive next problem.

\section{Why relative ML does not distinguish our curves}\label{sec:relativeML}
Let $A=k[x_1,\ldots,x_n]$, $I\subsetneq A$, and $\dim(A/I)\le n-2$. For every $r\ge1$,
\[
\boxed{\bigcap_{\substack{0\ne D\in\operatorname{LND}(A)\\D(A)\subset I^r}}\ker D=k.}
\]
Moreover, the algebra generated by the union of these kernels is $A$.

For each $i$, there exists a nonzero $f_i\in I\cap k[x_1,\ldots,\widehat{x_i},\ldots,x_n]$: otherwise, $n-1$ coordinate images would be algebraically independent in a ring of dimension $\le n-2$. Take
\[
D_i=f_i^r\partial_{x_i},\qquad
\ker D_i=k[x_1,\ldots,\widehat{x_i},\ldots,x_n].
\]
These LNDs satisfy $D_i(A)\subset I^r$; the intersection of their kernels is $k$, and their union contains every coordinate. This proves both assertions.

Consequently, the candidate with the weaker condition $D(I)\subset I$ also always gives $k$ for a curve in $\A^3$. Even requiring an arbitrary finite thickening to be fixed pointwise does not help.

For harmonic models, one may take, for example,
\[
f_x=C_c(y)-C_b(z),\quad
f_y=C_c(x)-C_a(z),\quad
f_z=C_b(x)-C_a(y).
\]
The mechanism of replicas is consistent with stronger results on flexibility of complements; see Flenner--Kaliman--Zaidenberg \cite{FKZ}. Thus the research strategy is corrected here; no new nontrivial invariant is being claimed.
